\scanpage{10}
\mspage{8}
une catégorie \uline{discrète}, ayant comme ens. d'\unsure{objets|objet} $\operatorname{Hom}_{X'}(x,y\unsure{\}|?})$. Ainsi le foncteur précédent définit en fait une application canonique
\[
\boxed{\pi_0\underline{\mathrm{Ch}}_\Sigma(X;x,y)\longrightarrow\operatorname{Hom}_\Sigma(x,y)}
\]
Je dis que cette application est toujours bijective. Pour le voir, nous allons \uline{construire} ``from scratch'' la catégorie $X\Sigma^{-1}$, en prenant comme objets ceux de $X$, et \deleted{\unsure{en posant}} en \uline{définissant}
\[
\boxed{\operatorname{Hom}_\Sigma(x,y)=\pi_0\underline{\mathrm{Ch}}_\Sigma(X;x,y)}
\]
Il faut donc définir la composition
\[
\operatorname{Hom}_\Sigma(x,y)\times\operatorname{Hom}_\Sigma(y,z)\longrightarrow\operatorname{Hom}_\Sigma(x,z).
\]

Pour ceci, on va définir des foncteurs ``composition des chemins''
\[
\begin{aligned}
&\deleted{\Phi_n\unsure{(}}\\
\underline{\mathrm{Ch}}_m(X;x,y)\times\underline{\mathrm{Ch}}_n(X;y,z)&\longrightarrow\underline{\mathrm{Ch}}_{\deleted{\unsure{m}}m+n}(X;x,z),\\
(c,\qquad c')&\longmapsto c'\circ c,
\end{aligned}
\]
mais \unsure{en|on} nous \unsure{bornant|bornons} pour simplifier au cas où $m,n$ ont la \uline{$=$ parité}, de telle façon qu'en ``mettant bout à bout'' $c$ et $c'$ \inserted{(dans cet ordre)}, on trouve bien un chemin de \unsure{type} $n+m$.

% archive sheet number: 9
